% !TeX root = ../main.tex

\chapter{数字问诊与临床训练}
\label{sec:digital-consultation}

普通学习平台擅长呈现教材、题目和讲解内容，却很少为学习者提供接近真实接诊节奏的连续问诊机会。学生即使已经记住疾病定义与治疗原则，也可能在面对一位只描述不适感受的患者时难以组织问题，难以从零散回答中提炼诊断依据，更缺少用完整语言说明鉴别诊断和治疗思路的反复练习。临床知识由书面记忆转向口头表达时，信息采集、推理判断和沟通措辞需要在同一过程内协同完成，这类能力很难通过单次选择题或静态病例阅读得到充分检验。

OralSage 为此设置独立的数字问诊页面，引入讯飞实时数字人扮演患者，学生可以使用语音或文字完成自由问诊。患者会围绕脱敏病例中的主诉、现病史和检查表现，用接近日常交流的语言回答追问，同时根据自然语言理解判断学生是否已经提出明确且正确的诊断。当诊断得到确认后，判官 Agent 会在原有对话中接续提问，引导学生说明病理机制、临床表现与诊断依据、鉴别诊断、治疗原则、风险与随访，使问诊信息采集和临床表达考核形成连贯体验。

这条训练链路将讯飞星火推理服务、流式语音听写和实时数字人能力组合到同一交互过程。系统在会话创建前检查病例、模型、语音和数字人服务的可用条件，交互数据经过结构校验与证据核验后才会推动流程进入下一阶段。每次诊断确认和分项作答都会形成专用学习事件，分数、遗漏要点与相关知识节点随之写入学习画像，后续选题可以重新读取这些变化，让一次临床表达训练自然影响下一轮病例安排。

\section{病例与个性化选题}

问诊病例来自 \texttt{data/processed/clinic\_stats/clinical\_cases.parsed.jsonl} 中的脱敏临床资料。病例进入训练池前需要同时具备明确诊断、主诉、现病史和检查记录，复诊随访材料与关键字段缺失的记录不会进入出题范围。当前训练池覆盖急性冠周炎、阻生牙、牙髓炎、急性根尖周炎、慢性根尖周炎、慢性牙周炎和根尖囊肿等常见口腔疾病，为主诉追问、病史采集和初步诊断提供了完整的文本依据。

会话启动时，选题服务读取学习画像中的 \texttt{known\_weak\_points}，并结合近期非问诊学习会话的主题形成学习目标。选题 Agent 将这些自然语言目标映射到训练池允许使用的疾病名称，随后由病例仓库从对应类别中选择一条符合训练条件的记录。学习画像暂时没有可用信息时，系统先随机选择疾病类别，再从该类别内抽取病例，使连续训练能够覆盖不同疾病，也避免记录数量较多的类别长期占据出题结果。

个性化选题与画像回写共同构成循环。一次训练暴露出的知识遗漏会进入薄弱点和知识节点掌握记录，后续会话再次读取这些信号，选题内容随学习表现逐步调整。病例对学生公开的内容只包含当前对话需要呈现的信息，目标诊断和内部证据保留在服务端上下文中，保证学习者需要通过问诊逐步完成判断。

\section{双路上下文构建}

患者扮演和临床评分需要不同类型的信息。系统在会话创建阶段并行准备两路上下文，一路从病例记录中整理患者能够陈述的事实，另一路围绕目标疾病从知识库检索教材和临床规范证据。两路信息分别交付给患者 Agent 与判官 Agent，病例叙述不会直接进入评分证据，教材知识也不会被拼接成患者自述，从数据边界上保护患者角色的真实感与评分依据的专业性。

\begin{figure}[H]
  \centering
  \begin{tikzpicture}[
    node distance=0.75cm and 1.1cm,
    every node/.style={font=\small,align=center},
    box/.style={draw=oralsagepurple,rounded corners=2pt,fill=coverpanel,minimum height=0.9cm,text width=3.0cm},
    agent/.style={draw=softcupblue,rounded corners=2pt,fill=softcupgray,minimum height=0.9cm,text width=2.8cm},
    line/.style={->,>=stealth,thick,color=covermuted}
  ]
    \node[box] (case) {脱敏病例\\主诉、病史与检查事实};
    \node[box,below=of case] (profile) {学习画像\\薄弱点与近期主题};
    \node[agent,right=of case] (patientctx) {患者上下文\\事实编号与对话历史};
    \node[agent,right=of profile] (judgectx) {判官上下文\\五个领域的检索证据};
    \node[box,right=of patientctx] (patient) {患者 Agent\\问诊与诊断理解};
    \node[box,right=of judgectx] (judge) {判官 Agent\\追问、评分与反馈};
    \draw[line] (case) -- (patientctx);
    \draw[line] (profile) -- (judgectx);
    \draw[line] (case) -- (judgectx);
    \draw[line] (patientctx) -- (patient);
    \draw[line] (judgectx) -- (judge);
    \draw[line] (patient) -- node[right,font=\scriptsize]{诊断确认} (judge);
  \end{tikzpicture}
  \caption{数字问诊的双路上下文构建}
  \label{fig:clinical-dual-context}
\end{figure}

患者上下文将病例字段转化为带编号的事实单元，并保留经过长度控制的会话历史。患者每次回答都要返回所依据的事实编号，事实核验 Agent 会检查回复能否由病例材料支持，经过核验的内容才会显示在对话中。遇到病例未提供的信息时，患者会以符合真实就诊交流的方式说明自己不清楚，避免补写病例中没有出现的检查结果或病程细节。

判官上下文围绕目标疾病分别检索病理机制、临床表现与诊断依据、鉴别诊断、治疗原则、风险与随访。检索结果经过疾病别名、相关度和来源类型筛选，临床病例类片段会被排除，至少三条有效教材或临床规范证据通过检查后才会建立考核题目。每条证据拥有编号和所属领域，后续题目、评分与批评意见均引用这些编号，使专业判断可以回溯到知识库资料。

\section{患者 Agent}

患者 Agent 使用讯飞星火 \texttt{xsparkx2flash} 推理服务理解学生的自然语言输入。Agent 会区分症状询问、诊断陈述、试探性诊断、鉴别讨论和其它交流意图，并以结构化字段返回动作、意图、诊断候选、诊断判断、患者回复、引用事实编号与未知信息标记。诊断是否正确由 Agent 对完整语义作出判断，学生可以使用“我考虑是……”“目前更倾向于……”等自然表达，无需记忆固定口令。

角色提示词要求患者始终使用第一人称，以生活化说法描述疼痛、肿胀、持续时间和诱发感受，并根据学生已经问到的内容控制回答范围。患者不会主动说出疾病名称，也不会对学生提出的诊断作确认式回应。即使学生直接询问某个疾病名称，回复仍会停留在患者可感知的症状和经历中，诊断判断只通过内部字段交给状态机处理。

每轮对话同时检查事实依据与诊断判断。事实核验 Agent 关注患者回复中出现的陈述能否映射到病例事实，诊断复核 Agent 关注学生是否已经给出明确诊断以及候选疾病是否与病例目标一致。两项结果经过结构检查后合并到当前轮次，既保留自然交流的流畅感，也让诊断交接建立在完整语义和病例内容之上。

\section{判官 Agent}

学生完成诊断后，判官 Agent 沿用当前对话页面接管后续交流。系统根据检索证据预先构建五道考核问题，依次覆盖病理机制、临床表现与诊断依据、鉴别诊断、治疗原则、风险与随访，每个领域满分 20 分。题目包含唯一编号、领域名称、提问文本、评分要点、证据引用和最高分，五个领域共同形成 100 分的临床表达评价。

判官对每次回答返回数值分数、针对性评语、已经表达清楚的内容、尚未覆盖的要点、需要修正的表述和证据编号。评分引用只能来自当前题目允许使用的证据集合，服务端会检查分数范围、问题顺序和引用合法性。学生能够在对话中立即看到与自己表述相对应的反馈，随后的问题继续承接同一疾病语境，完成全部五项后生成总分和分项结果。

每个分项结果都会更新相关知识节点的掌握度、接触次数和置信度，遗漏要点进入已知薄弱点。诊断结果、领域分数、总分和批评意见也会作为专用临床训练事件出现在学习画像时间线中。判官输出由此兼顾当次表达指导和长期学习记录，使学生在下一次课程复习或病例训练中能够继续处理已经识别出的知识缺口。

\section{状态机设计}

数字问诊服务采用四阶段状态机管理患者交互与考核交接，合法状态依次为 \texttt{consultation}、\texttt{diagnosis\_confirmation}、\texttt{oral\_exam} 和 \texttt{completed}。患者 Agent 在问诊阶段持续回应学生追问，识别到明确且正确的诊断后进入诊断确认，学生确认诊断便触发判官 Agent 建立五域考核，五道题完成评分后进入完成状态。

\begin{center}
  \textbf{问诊交互 $\rightarrow$ 诊断识别 $\rightarrow$ 用户确认 $\rightarrow$ 五域表达考核 $\rightarrow$ 画像更新}
\end{center}

会话状态保存轮次记录、当前问题序号、历次评分、诊断候选和已处理请求编号。服务端只接受与当前阶段相符的操作，问题按既定顺序推进，已经完成的请求再次到达时会返回同一结果。SQLite 持久化使刷新页面或重新打开历史记录后仍可恢复患者对话、考核进度和评分信息，数字人播报只消费已经进入会话记录的 Agent 回复。

状态变化还对应四类学习事件：\texttt{clinical\_speaking\_started} 记录训练启动，\texttt{clinical\_diagnosis\_confirmed} 记录诊断确认，\texttt{clinical\_speaking\_answer\_assessed} 保存分项评分，\texttt{clinical\_speaking\_completed} 汇总训练结果。独立事件类型让学习画像可以按“问诊、诊断、考核”呈现活动，也为后续个性化选题提供清晰的数据来源。

\section{讯飞能力集成}

数字问诊集中使用讯飞开放平台的推理、语音和数字人产品。患者 Agent、诊断复核、选题映射、判官题目构建与评分均调用星火 \texttt{xsparkx2flash}，该模型对中文自然表达和多轮上下文的理解使学生可以按照真实问诊习惯组织语言。服务通过 \texttt{xfyun\_spark\_reasoning\_http} 适配器接入推理接口，专用环境变量与课程学习模型配置分离，数字问诊可以独立校验模型名称、服务地址和访问密钥。

学生选择语音输入时，浏览器连接讯飞流式语音听写 WebSocket，普通话音频会在说话过程中连续转写，停顿检测参数适配较长的医学问句。服务端生成短时有效的签名地址，AppID、APIKey 和 APISecret 不会进入页面配置。转写文本提交给患者 Agent 前还会经过医学术语整理，牙体编号、疾病名称和口腔专有词优先保留，减少同音词对诊断语义的影响。

实时数字人 SDK 在浏览器中建立 \texttt{xrtc} 会话，数字人使用 \texttt{x4\_lingxiaoqi\_oral} 发音人播报患者或判官回复。患者形象采用透明渲染层，与页面的淡紫色视觉体系自然融合，学生可以在同一视线范围内观察数字人表达并继续输入问题。星火生成的生活化回答随会话记录交给数字人播报，语音听写又把学生口述送回 Agent，上下行能力结合后形成接近面对面交流的训练节奏。

\begin{table}[H]
  \centering
  \caption{讯飞能力在数字问诊中的作用}
  \label{tab:clinical-xfyun-capabilities}
  \begin{tabularx}{\linewidth}{p{0.27\linewidth}p{0.28\linewidth}X}
    \toprule
    讯飞产品 & 接入方式 & 训练效果 \\
    \midrule
    星火 \texttt{xsparkx2flash} & 推理服务 HTTP 接口 & 理解多轮问诊、自然语言诊断与临床表达，并生成受结构约束的患者和判官结果 \\
    流式语音听写 & \texttt{wss://iat-api.xfyun.cn/v2/iat} & 将连续口述转化为可编辑文本，医学术语整理后进入问诊会话 \\
    实时数字人 & 实时互动 WebSocket 与浏览器 SDK & 以患者形象播报 Agent 回复，透明画面与问诊页面保持一致 \\
    长文本语音合成 & 异步语音合成接口 & 为课程讲解生成可反复收听的音频，与数字问诊共同覆盖听说训练和知识学习 \\
    \bottomrule
  \end{tabularx}
\end{table}

\section{稳定性设计}

数字问诊在创建会话前执行完整的就绪检查，检查内容包括星火推理服务授权、语音听写配置、实时数字人配置和可训练病例数量。页面只在必要条件满足后开放新会话，配置问题会以明确字段返回，便于部署人员直接定位对应的讯飞服务或本地数据项。

Agent 之间使用严格 JSON 契约传递选题、患者回答、诊断交接、考核题目和评分结果。服务端会逐项校验字段类型、枚举值、事实编号、证据引用、分数范围和问题顺序，只有通过检查的数据才会写入会话并推动状态变化。上下文保留范围和模型输出长度均设置明确上限，检索证据按领域筛选后再交付给判官，减少冗余信息对推理结果的干扰。

每次创建会话、提交问诊轮次和提交考核回答都携带请求编号，同一会话的写操作由操作锁串行处理。重复请求会读取已保存的处理结果，学习画像写入使用可复现的事件编号，网络重发不会造成轮次、分数或掌握记录重复。会话与判官题目蓝图保存在 \texttt{clinical\_speaking.sqlite3}，相同疾病和证据签名可以复用已校验的五域题目，页面刷新后也能继续当前阶段。

对讯飞推理服务的短暂网络异常采用次数受限的重试，并对空响应和不完整结构再次请求有效结果。每次请求的最终内容仍需经过原有结构与证据检查，未通过检查的数据不会进入对话记录。数字问诊专项测试共 49 项，覆盖自然语言诊断理解、患者角色边界、事实核验、双路检索、患者与判官交接、五域评分、语音术语、重复请求、状态恢复、画像写入和数字人鉴权，验证了这条训练链路的关键行为。
